% Source: https://artofproblemsolving.com/community/c3h1573920_advanced_hard_algebriac_fractions
% Stable identity: aops-topic-1573920-q01
\textbf{Clearing denominators.} The equation only makes sense for $n\ne1$ and $n\ne2$. For such $n$ the number $(n-1)(n-2)$ is nonzero, and multiplying both sides of an equation by the same nonzero number gives an equivalent equation. We choose this number because it is a common multiple of both denominators, so every fraction disappears:
\[
\dfrac{9n}{n-2}\cdot(n-1)(n-2)=9n(n-1),\qquad \dfrac{8n}{n-1}\cdot(n-1)(n-2)=8n(n-2),\qquad 1\cdot(n-1)(n-2)=n^2-3n+2 .
\]
The term $n^2-3n+2$ is where the ``$+1$'' goes: it must also be multiplied by $(n-1)(n-2)$.

\textbf{Solving.}
\[
9n^2-9n=8n^2-16n+n^2-3n+2\quad\Longrightarrow\quad -9n=-19n+2\quad\Longrightarrow\quad 10n=2,
\]
so $n=\boxed{\dfrac{1}{5}}$. This is not one of the excluded values $1$, $2$, so it is a genuine solution.

Check: $\dfrac{9n}{n-2}=\dfrac{\dfrac{9}{5}}{-\dfrac{9}{5}}=-1$ and $\dfrac{8n}{n-1}+1=\dfrac{\dfrac{8}{5}}{-\dfrac{4}{5}}+1=-2+1=-1$.
